% preâmbulo comum — ambiente sem lmodern: ae + aecompl (regras do projeto)
\documentclass[11pt,a4paper]{article}
\usepackage{ae,aecompl}
\usepackage[utf8]{inputenc}
\IfFileExists{brazil.ldf}{\usepackage[brazil]{babel}}{}
\usepackage[a4paper,margin=2.3cm]{geometry}
\usepackage{amsmath,amssymb,amsthm,mathtools}
\usepackage{booktabs,array,longtable}
\usepackage{enumitem}
\usepackage{xcolor}
\usepackage{graphicx}
\usepackage{tikz}
\usetikzlibrary{arrows.meta,positioning,fit,calc}
\usepackage[most]{tcolorbox}
\usepackage{listings}
\usepackage[ruled,vlined,linesnumbered]{algorithm2e}
\usepackage{microtype}
\usepackage{fancyhdr}
\usepackage{float}
\usepackage[colorlinks=true,linkcolor=blue!60!black,citecolor=blue!60!black,urlcolor=blue!60!black]{hyperref}
\setlength{\emergencystretch}{1.5em}
% marcadores em modo matemático (TS1 sem face vetorial sob ae)
\renewcommand{\labelitemi}{$\bullet$}
\renewcommand{\labelitemii}{$\circ$}
\renewcommand{\labelitemiii}{$\ast$}
% nomes em português sem depender do babel
\renewcommand{\contentsname}{Sumário}
\renewcommand{\refname}{Referências}
\renewcommand{\figurename}{Figura}
\renewcommand{\tablename}{Tabela}
\renewcommand{\abstractname}{Resumo}
\renewcommand{\proofname}{Demonstração}
\SetAlgorithmName{Algoritmo}{Algoritmo}{Lista de Algoritmos}
\definecolor{tune3blue}{RGB}{31,73,125}
\definecolor{okgreen}{RGB}{20,110,60}
\definecolor{warnred}{RGB}{160,30,30}
\definecolor{codebg}{RGB}{245,246,248}
\lstset{basicstyle=\ttfamily\small,backgroundcolor=\color{codebg},frame=single,rulecolor=\color{gray!40},
  breaklines=true,columns=fullflexible,keepspaces=true,showstringspaces=false,
  language=Python,keywordstyle=\color{tune3blue}\bfseries,commentstyle=\color{gray!70!black}\itshape,
  literate={á}{{\'a}}1 {à}{{\`a}}1 {ã}{{\~a}}1 {â}{{\^a}}1 {é}{{\'e}}1 {ê}{{\^e}}1 {í}{{\'i}}1 {ó}{{\'o}}1 {ô}{{\^o}}1 {õ}{{\~o}}1 {ú}{{\'u}}1 {ç}{{\c{c}}}1 {Á}{{\'A}}1 {É}{{\'E}}1 {Í}{{\'I}}1 {Ó}{{\'O}}1 {Ú}{{\'U}}1 {Ç}{{\c{C}}}1 {Ã}{{\~A}}1 {Õ}{{\~O}}1 {Ê}{{\^E}}1 {Ô}{{\^O}}1 {Â}{{\^A}}1 {²}{{$^2$}}1 {·}{{$\cdot$}}1 {≥}{{$\geq$}}1 {≤}{{$\leq$}}1 {→}{{$\to$}}1 {∈}{{$\in$}}1 {∞}{{$\infty$}}1 {σ}{{$\sigma$}}1 {ρ}{{$\rho$}}1 {γ}{{$\gamma$}}1 {η}{{$\eta$}}1 {θ}{{$\theta$}}1 {Σ}{{$\Sigma$}}1 {λ}{{$\lambda$}}1 {ε}{{$\varepsilon$}}1 {–}{{--}}1 {—}{{---}}1 {×}{{$\times$}}1}
\newtcolorbox{ideia}[1][]{colback=tune3blue!5,colframe=tune3blue,fonttitle=\bfseries,title={#1},breakable}
\newtcolorbox{alerta}[1][]{colback=warnred!5,colframe=warnred,fonttitle=\bfseries,title={#1},breakable}
\newtcolorbox{codigo}[1][]{colback=okgreen!4,colframe=okgreen,fonttitle=\bfseries,title={#1},breakable}
\theoremstyle{definition}
\newtheorem{definicao}{Definição}[section]
\newtheorem{proposicao}{Proposição}[section]
\newtheorem{exemplo}{Exemplo}[section]
\newtheorem{observacao}{Observação}[section]
\DeclareMathOperator{\Tr}{Tr}
\DeclareMathOperator{\CVaR}{CVaR}
\DeclareMathOperator{\VaR}{VaR}
\DeclareMathOperator{\Var}{Var}
\DeclareMathOperator{\Cov}{Cov}
\DeclareMathOperator{\GSNR}{GSNR}
\DeclareMathOperator*{\argmin}{arg\,min}
\newcommand{\R}{\mathbb{R}}
\newcommand{\E}{\mathbb{E}}
\newcommand{\norm}[1]{\left\lVert #1 \right\rVert}
\newcommand{\code}[1]{\texttt{#1}}
\pagestyle{fancy}\fancyhf{}
\rhead{\small Tune3 — \docshort}
\lhead{\small UNIPAMPA / AI Horizon Labs}
\cfoot{\thepage}
